% Plantilla del paper final (IEEE, 4 a 6 páginas). SOLO ESTRUCTURA.
% El texto lo redacta el grupo: las políticas del curso prohíben la
% redacción automatizada del paper final. Cada sección lleva, como
% comentario, quién la redacta, quién la revisa, de dónde sale y qué
% debe cubrir, según la "Nueva distribución del trabajo y del paper".
%
% Compilar en Overleaf (IEEEtran ya viene instalado): pdfLaTeX + BibTeX.
% Las tablas de tablas/ y las figuras de figuras/ salen del notebook:
%   trabajo/resultados/<particion>/tablas/*.tex  ->  tablas/
%   trabajo/resultados/<particion>/fig_*.png      ->  figuras/
%   trabajo/resultados/fig_prefijos.png           ->  figuras/
% Mientras falte un archivo, se dibuja un recuadro con su nombre.
%
% Cada sección está en su propio archivo de secciones/ (el comentario de arriba
% de cada archivo dice quién la redacta y quién la revisa). Así cada integrante
% edita solo su archivo y nadie pisa el trabajo de otro.
%
% Numeración: LaTeX numera figuras y tablas en el orden en que se DEFINEN en
% este archivo; por eso los pies no llevan número y en el texto se cita con
% \ref{...}. El orden de abajo da la numeración del plan: Fig. 1 prefijos,
% 2 curvas, 3 por k, 4 t-SNE, 5 H4, 6 H5, 7 trazo a trazo; Tabla I búsqueda,
% II modelos, III protocolo, IV hipótesis.

\documentclass[conference]{IEEEtran}
\usepackage[utf8]{inputenc}
\usepackage[T1]{fontenc}
\usepackage[spanish,es-nodecimaldot,es-noshorthands]{babel}
\usepackage{amsmath,amssymb}
\usepackage{graphicx}
\usepackage{booktabs}
\usepackage{url}
\usepackage{xcolor}

% \figura{archivo}{ancho}: la imagen si existe; si no, un recuadro con el nombre.
\newcommand{\figura}[2]{%
  \IfFileExists{figuras/#1}{\includegraphics[width=#2]{figuras/#1}}{%
    \fbox{\parbox[c][3cm][c]{\dimexpr#2-2\fboxsep-2\fboxrule\relax}{\centering\small\texttt{\detokenize{#1}}\\falta la figura}}}}
% \tabla{archivo}: la tabla si existe; si no, un aviso.
\newcommand{\tabla}[1]{%
  \IfFileExists{tablas/#1}{\input{tablas/#1}}{\fbox{\small\texttt{\detokenize{#1}}: falta la tabla}}}
% \pendiente{texto}: marca en rojo lo que falta (borrar antes de entregar).
\newcommand{\pendiente}[1]{\textcolor{red}{[#1]}}

\begin{document}

\title{¿Cuántos trazos bastan? Reconocimiento temprano de bocetos con CNN en Quick, Draw!}

\author{
\IEEEauthorblockN{Elias Enoc Rosales Torres, Jorly Inga Chavez, Farid Alonzo Villanueva Alzamora,\\
Cristian Brayan Diaz Pariacuri y Omar Ailton Ccahuana Sosa}
\IEEEauthorblockA{Escuela Profesional de Ciencia de la Computación, Facultad de Ciencias\\
Universidad Nacional de Ingeniería, Lima, Perú\\
Inteligencia Artificial (CC421), 2026-II}
}

\maketitle

% Resumen · Redacta: Farid · Revisan: todos · Semana 14 · 150 a 200 palabras.
% Debe cubrir: la pregunta, los datos (30 categorías), los tres modelos,
% el resultado principal de cada hipótesis con su número, la conclusión.
\begin{abstract}
\pendiente{Resumen}
\end{abstract}

\begin{IEEEkeywords}
\pendiente{Palabras clave}
\end{IEEEkeywords}

% I · Redacta: Elias · Revisa: Jorly · Semana 9 · ~0,4 pág.
% Sale de: avance 1.1 a 1.3 y 1.6.
% Debe cubrir: el juego y su mecánica (la ronda termina cuando la IA
% acierta), el problema en sus tres preguntas, la pregunta de investigación,
% las contribuciones (de los objetivos 1.5, sin M0 ni PCA) y la
% organización del paper.
\section{Introducción}
\pendiente{Introducción}

% II · Redacta: Jorly · Revisa: Farid · Semana 9 · ~0,3 pág.
% Sale de: avance 2.1, 2.2 y 2.4. Sin la calibración (2.3).
% Debe cubrir: Sketch-a-Net, SketchMamba y en qué se diferencia este trabajo.
\section{Trabajo relacionado}
\pendiente{Trabajo relacionado}

% III · Metodología. La teoría necesaria va dentro de cada subsección.
\section{Metodología}

% III-A · Redacta: Elias · Revisa: Omar · Semana 9 · ~0,3 pág.
% Sale de: avance 3.1 y 3.6.1; docs/01_datos.md.
% Debe cubrir: el dataset y sus campos, el criterio de selección de las 30
% categorías (langosta por trompeta, violín como excepción), el muestreo
% por reservorio, las particiones 70/15/15 por dibujo y la tabla de hechos
% resumida.
\subsection{Datos}
\pendiente{Datos}

% III-B · Redacta: Jorly · Revisa: Cristian · Semana 9 · ~0,3 pág.
% Sale de: avance 3.2 (ecuaciones 1 a 3); docs/02_prefijos_y_rasterizacion.md.
% Debe cubrir: prefijo D_k, normalización causal, rasterización a 28x28
% (lienzo 256, grosor 16, margen 2; el avance decía 1, ver docs/02) y la Fig. 1.
\subsection{Prefijos y normalización causal}
\pendiente{Prefijos y normalización causal}

\begin{figure}[t]
\centering
\figura{fig_prefijos.png}{\columnwidth}
\caption{\pendiente{Prefijos con normalización causal (fila superior) frente a la caja del dibujo completo.}}
\label{fig:prefijos}
\end{figure}

% III-C · Redacta: Cristian · Revisa: Elias · Semana 10 · ~0,4 pág.
% Sale de: avance 3.3 y 3.4 (Cuadro 3); docs/03 y docs/04.
% Debe cubrir: softmax y entropía cruzada en una línea, M1, la CNN de M2
% y M2-P (Cuadro 3, GAP), el esquema mixto de M2-P y por qué se cambió desde
% el de solo prefijos (regla fijada antes, docs/04 "Decisión sobre H2"), la búsqueda de
% hiperparámetros, las semillas, la parada temprana y UNA línea que diga
% que el alcance se redujo respecto al avance y por qué.
\subsection{Modelos y entrenamiento}
\pendiente{Modelos y entrenamiento}

% III-D · Redacta: Omar · Revisa: Farid · Semana 10 · ~0,4 pág.
% Sale de: avance 3.6.2, 3.6.3 y 3.7 (ecuaciones 9 y 10, Cuadro 4); docs/05 y docs/06.
% Debe cubrir: exactitud por fracción, top-3, AUC progresiva, k*, bootstrap,
% McNemar y los criterios de H1, H2 y H5 (en las tres semillas).
\subsection{Métricas y pruebas estadísticas}
\pendiente{Métricas y pruebas estadísticas}

% III-E · Redacta: Farid · Revisa: Omar · Semana 10 · ~0,2 pág.
% Sale de: plan v5 ("H4 reformulada"); docs/07_analisis_h4.md.
% Debe cubrir: confundibilidad con kNN en el espacio GAP de M2, mediana de
% k*/K de M2-P, Spearman, por qué no es circular y t-SNE solo para mirar.
\subsection{Confundibilidad en el espacio de la CNN}
\pendiente{Confundibilidad}

% IV · Resultados. Hasta la semana 12 con números de VALIDACIÓN;
% en la semana 13 se reemplazan por los de PRUEBA (regenerar tablas/ y figuras/).
\section{Resultados}

% IV-A · Redacta: Cristian · Revisa: Omar · Semana 11 · ~0,2 pág.
\subsection{Búsqueda de hiperparámetros}
\pendiente{Búsqueda de hiperparámetros}

\begin{table}[t]
\centering
\caption{\pendiente{Búsqueda de hiperparámetros en validación.}}
\label{tab:busqueda}
\resizebox{\columnwidth}{!}{\tabla{tabla_busqueda.tex}}
\end{table}

% IV-B · Redacta: Omar · Revisa: Cristian · Semanas 12 y 13 · ~0,5 pág.
% Debe cubrir: Tabla 2, Figs. 2, 3 y 6, y el veredicto de H1, H2 y H5 con su número.
\subsection{¿Con cuántos trazos? (H1, H2 y H5)}
\pendiente{H1, H2 y H5}

\begin{table}[t]
\centering
\caption{\pendiente{M1, M2 y M2-P, media $\pm$ desviación de las semillas.}}
\label{tab:modelos}
\resizebox{\columnwidth}{!}{\tabla{tabla_modelos.tex}}
\end{table}

\begin{figure}[t]
\centering
\figura{fig_curvas_fraccion.png}{\columnwidth}
\caption{\pendiente{Exactitud por fracción del dibujo (figura principal).}}
\label{fig:curvas}
\end{figure}

\begin{figure}[t]
\centering
\figura{fig_exactitud_por_k.png}{\columnwidth}
\caption{\pendiente{Exactitud por número de trazos.}}
\label{fig:k}
\end{figure}

% IV-C · Redacta: Farid · Revisa: Jorly · Semanas 12 y 13 · ~0,3 pág.
% Debe cubrir: Figs. 4 y 5, el veredicto de H4 y la separación de pares
% (separacion_pares.csv).
\subsection{¿Por qué algunas categorías necesitan más? (H4)}
\pendiente{H4 y análisis por par}

\begin{figure*}[t]
\centering
\figura{fig_tsne.png}{0.85\textwidth}
\caption{\pendiente{t-SNE de los vectores GAP de M2-P al 30\,\% y al 100\,\% del dibujo.}}
\label{fig:tsne}
\end{figure*}

\begin{figure}[t]
\centering
\figura{fig_h4.png}{\columnwidth}
\caption{\pendiente{H4: confundibilidad frente a la mediana de $k^*/K$.}}
\label{fig:h4}
\end{figure}

% Fig. 6 (H5) y Fig. 7 (trazo a trazo) se definen aquí, después de H4, para
% que la numeración siga el plan; en IV-B se citan con \ref{fig:h5}.
\begin{figure}[t]
\centering
\figura{fig_h5.png}{\columnwidth}
\caption{\pendiente{H5: distribución de $K-k^*$ (trazos de anticipación).}}
\label{fig:h5}
\end{figure}

% Fig. 7 · Prepara: Jorly. Predicción de M1, M2 y M2-P en cada trazo de
% cuatro dibujos de los pares; respaldo de la demo en la exposición.
\begin{figure*}[t]
\centering
\figura{fig_trazo_a_trazo.png}{0.85\textwidth}
\caption{\pendiente{Predicción de cada modelo trazo a trazo en dibujos de los pares.}}
\label{fig:trazo}
\end{figure*}

% V · Redactan: todos, cada uno sobre su hipótesis o módulo; Elias junta
% las limitaciones · Revisa: Farid · Semana 14 · ~0,4 pág.
% Sale de: avance 1.6 y los criterios de cada hipótesis.
% Debe cubrir: qué significa cada veredicto, la Tabla 3 (protocolo frente
% a SketchMamba), las limitaciones (país, recognized, dibujos cortados,
% simplificación, H4 depende de una CNN, 3 semillas) y trabajo futuro.
\section{Discusión y limitaciones}
\pendiente{Discusión y limitaciones}

\begin{table}[t]
\centering
\caption{\pendiente{Diferencias de protocolo con SketchMamba (prepara Jorly).}}
\label{tab:protocolo}
\pendiente{Tabla 3}
\end{table}

\begin{table}[t]
\centering
\caption{\pendiente{Resumen de las hipótesis por semilla.}}
\label{tab:hipotesis}
\resizebox{\columnwidth}{!}{\tabla{tabla_hipotesis.tex}}
\end{table}

% VI · Redacta: Farid · Revisan: todos · Semana 14 · ~0,2 pág.
\section{Conclusiones}
\pendiente{Conclusiones}

% Opcional, según lo que pida el curso. La DUIA va aparte con firmas (Nivel 3).
% \section*{Declaración de uso de IA}

\bibliographystyle{IEEEtran}
\bibliography{referencias}

\end{document}
